\section{Parité}
Les tests que nous avons faits jusqu'à présent dans les instructions
``if'' concernaient l'orientation du \robot.  En \ocaml\ il est
possible de faire des tests concernant les valeurs des paramètres et des
compteurs de boucles ``for''.  Il peut être utile de faire un test sur le compteur
dans le corps d'une boucle pour exécuter des actions différentes à
chaque itération.  Prenons l'exemple du dessin de la frise en forme de
double peigne ci-dessous~:

 {\raggedright
  \includevspacePostScript[1.5cm]{FriseDoublePeigne}{1.0}{1.0}
   }

Les dents du peigne étant dessinées de gauche à droite, les
$1^{\mathrm{\grave{e}re}}$, $3^{\mathrm{\grave{e}me}}$,
$5^{\mathrm{\grave{e}me}}$, \ldots\ sont dirigées vers le bas alors
que les $2^{\mathrm{\grave{e}me}}$, $4^{\mathrm{\grave{e}me}}$,
$6^{\mathrm{\grave{e}me}}$, \ldots\ sont dirigées vers le haut.  Dans
le programme de dessin on utilise un compteur de boucle {\idt i} qui
vaut 1 quand on dessine la $1^{\mathrm{\grave{e}re}}$ dent, 2 quand on
dessine la $2^{\mathrm{\grave{e}me}}$, \ldots\ Par conséquent quand la
valeur du compteur de boucle est {\em impaire} (c'est-à-dire 1, 3, 5,
\ldots) il faut dessiner une dent vers le bas tandis que si la valeur
du compteur est {\em paire} (c'est-à-dire 2, 4, 6, \ldots) il faut
dessiner une dent vers le haut.  \index{Parité}Un nombre est impair
quand le reste de sa division par 2 est égal à 1. Pour tester que la
valeur du compteur {\idt i} est impaire on utilise donc l'expression
booléenne {\idt i mod 2 = 1} qui vaut {\idt true} (vrai) quand le reste
{\idt i mod 2} de la division entière de {\idt i} par 2 est égal à 1 et
{\idt false} (faux) quand ce reste est égal à 0, c'est-à-dire quand
{\idt i} est pair.
\newpage

{
\parindent \parindentprog
\medskip

\carrenoir
\vspace{\neglskip}\mbox{\blc {\mcl open}{\idt \ LibRobot}}\par
\vspace{0.3\baselineskip}
\vspace{\neglskip}\mbox{\blc {\cmt (*\ Dessiner\ une\ frise\ en\ forme\ de\ double\ peigne\ *)}}\par
\vspace{0.3\baselineskip}
\vspace{\neglskip}\mbox{\blc {\mcl let}{\idt \ nombreDeDents\ =\ 30}}\par
\vspace{0.3\baselineskip}
\vspace{\neglskip}\mbox{\blc {\mcl let}{\idt \ ()\ =}}\par
\vspace{\neglskip}\mbox{\blc \blk{\idt ag\ ();\ pe\ ();}}\par
\vspace{\neglskip}\mbox{\blc \blk{\mcl for}{\idt \ i\ =\ 1\ }{\mcl to}{\idt \ nombreDeDents\ }{\mcl do}}\par
\vspace{\neglskip}\mbox{\blc \blk\blk{\cmt (*\ Trait\ horizontal\ *)}}\par
\vspace{\neglskip}\mbox{\blc \blk\blk{\idt av\ ();}}\par
\vspace{\neglskip}\mbox{\blc \blk\blk{\mcl if}{\idt \ i\ }{\mcl mod}{\idt \ 2\ =\ 1\ }{\mcl then}{\idt \ }{\mcl begin}}\par
\vspace{\neglskip}\mbox{\blc \blk\blk\blk{\cmt (*\ Dent\ vers\ le\ bas\ *)}}\par
\vspace{\neglskip}\mbox{\blc \blk\blk\blk{\idt pqtd\ ();\ av\ ();\ pdt\ ();\ av\ ();\ pqtd\ ()}}\par
\vspace{\neglskip}\mbox{\blc \blk\blk{\mcl end}{\idt \ }{\mcl else}{\idt \ }{\mcl begin}}\par
\vspace{\neglskip}\mbox{\blc \blk\blk\blk{\cmt (*\ Dent\ vers\ le\ haut\ *)}}\par
\vspace{\neglskip}\mbox{\blc \blk\blk\blk{\idt pqtg\ ();\ av\ ();\ pdt\ ();\ av\ ();\ pqtg\ ()}}\par
\vspace{\neglskip}\mbox{\blc \blk\blk{\mcl end}}\par
\vspace{\neglskip}\mbox{\blc \blk{\mcl done}{\idt ;}}\par
\vspace{\neglskip}\mbox{\blc \blk{\idt st\ ()}}\par

\medskip }


\section{Comparaison d'entiers}

Si {\idt i} est un compteur de boucle ou un paramètre entier de procédure,
l'expression booléenne {\idt i > 10} est {\idt true} (vraie)
si la valeur de {\idt i} est strictement plus grande que 10 et {\idt
false} (fausse) dans le cas contraire c'est-à-dire si la valeur de {\idt
i} est inférieure ou égale à 10.  Les symboles mathématiques $\leq$
(inférieur ou égal) et $\geq$ (supérieur ou égal) se notent {\idt <=} et
{\idt >=} en \ocaml\ (car les symboles mathématiques $\leq$ et $\geq$ ne
se trouvent pas sur les claviers d'ordinateurs alors que {\idt <}, {\idt >}
et {\idt =} sont toujours présents).  Par exemple les expressions
booléennes {\idt 10 <= 10} et {\idt 10 <= 20} valent {\idt true}
(vrai) tandis que {\idt 10 <= 9} vaut {\idt false} (faux).  En
\ocaml\ on note l'égalité avec le symbole mathématique usuel {\idt =}.
Par contre le symbole mathématique $\neq$ (différent) se note {\idt <>}.
Par exemple {\idt 10 = 10} et {\idt 10 <> 20} valent {\idt true}
(vrai) tandis que {\idt 10 <> 10} et {\idt 10 = 20} valent {\idt
false} (faux)\footnote{\ On rencontre aussi en \ocaml\ les symboles
{\idt ==} et {\ttfamily\idt !=}, qui ont un autre sens. Pour comparer des
entiers, il faut toujours utiliser {\idt =} et {\idt <>}.}.
On peut utiliser les expressions booléennes de
comparaison d'entiers ou de valeurs de paramètres entiers dans le test
de l'instruction ``if''.  Par exemple la commande {\idt avf n} qui
fait avancer le \robot\ {\idt n} fois si la valeur de {\idt n} est
positive, le fait reculer si la valeur de {\idt n} est négative et qui
ne change pas sa position quand la valeur de {\idt n} est nulle peut
se définir comme suit~:

\input ocaml/DefinitionAvf.tex

\noindent On remarquera que le paramètre effectif {\idt -5} de l'appel
{\idt avf (-5)} doit être placé entre parenthèses. Si l'on écrit
{\idt avf -5}, le compilateur comprend que l'on veut soustraire 5 de
{\idt avf}$\:$!

De manière plus générale, on peut utiliser les opérateurs {\idt <} (strictement inférieur), {\idt <=} (inférieur ou égal),
{\idt =} (égal), {\idt <>} (différent), {\idt >=} (supérieur ou égal) et {\idt >} (strictement supérieur) pour comparer les valeurs
de deux expressions entières comme par exemple {\idt i <= l} dans le programme suivant qui trace le serpentin
ci-dessous. La commande {\idt lg} attend un nombre écrit avec un point, comme
dans {\idt lg 20.}$\:$; pour lui donner la valeur d'une expression entière comme
{\idt i}, on la convertit avec {\idt float\_of\_int} (comme nous l'expliquerons
au chapitre 19, consacré aux expressions rationnelles)~:

 \includevspacePostScript[1.5cm]{Serpentin}{0.95}{0.95}
 \input ocaml/Serpentin.tex

\section{Négation, conjonction et disjonction}

La {\em négation} \index{Négation} d'une expression booléenne $E$ s'écrit {\mcl not} $E$, le mot anglais {\mcl not} signifiant `non'. Si la
valeur de $E$ est {\idt true} alors celle de {\mcl not} $E$ est {\idt false} tandis que si la valeur de $E$ est {\idt false} alors
celle de {\mcl not} $E$ est {\idt true}. Par exemple  {\mcl not} {\idt (5 mod 2 = 1)} vaut {\idt false} puisque
{\idt 5 mod 2 = 1} vaut {\idt true}, {\idt 5} étant impair. Les parenthèses sont
indispensables (nous verrons pourquoi dans l'Addenda). Ceci peut se résumer par la table de vérité suivante~:

\begin{center}
\begin{tabular}{|c|c|} \hline
$E$ & {\mcl not} $E$\\ \hline \hline
{\idt true} & {\idt false}\\ \hline
{\idt false} & {\idt true}\\ \hline
\end{tabular}
\end{center}

La {\em conjonction} \index{Conjonction} de deux expressions booléennes $E_{1}$ et $E_{2}$ s'écrit $E_{1}$ {\idt \&\&} $E_{2}$,
ce qui se lit `$E_{1}$ et $E_{2}$'. La valeur de $E_{1}$ {\idt \&\&} $E_{2}$ est {\idt true} (vrai) si et
seulement si la valeur de $E_{1}$ et celle de $E_{2}$ sont {\idt true} (vrai). Si l'une des expressions
booléennes $E_{1}$ ou
$E_{2}$ vaut {\idt false} (faux) alors la conjonction $E_{1}$ {\idt \&\&} $E_{2}$ vaut également {\idt false} (faux).
Par exemple  {\idt 5 >= 0 \&\& 5 mod 2 = 1} vaut {\idt true} puisque
{\idt 5} est positif et impair. Ceci peut se résumer par la table de vérité suivante~:

\begin{center}
\begin{tabular}{|c|c|c|} \hline
$E_{1}$ & $E_{2}$ & $E_{1}$ {\idt \&\&} $E_{2}$\\ \hline \hline
{\idt true} & {\idt true} & {\idt true}\\ \hline
{\idt true} & {\idt false} & {\idt false}\\ \hline
{\idt false} & {\idt true} & {\idt false}\\ \hline
{\idt false} & {\idt false} & {\idt false}\\ \hline
\end{tabular}
\end{center}

La {\em disjonction} \index{Disjonction} de deux expressions booléennes $E_{1}$ et $E_{2}$ s'écrit $E_{1}$ {\idt ||} $E_{2}$,
ce qui se lit `$E_{1}$ ou $E_{2}$'. La valeur de $E_{1}$ {\idt ||} $E_{2}$ est {\idt true} (vrai) si et
seulement si la valeur de $E_{1}$ ou celle de $E_{2}$ est {\idt true} (vrai). Si les valeurs de $E_{1}$ et de
$E_{2}$ sont  {\idt false} (faux) alors la disjonction $E_{1}$ {\idt ||} $E_{2}$ vaut également {\idt false} (faux).
Par exemple  {\idt 4 >= 0 || 4 mod 2 = 1} vaut {\idt true} puisque
{\idt 4} est positif et bien que {\idt 4} soit pair. Ceci peut se résumer par la table de vérité suivante~:

\begin{center}
\begin{tabular}{|c|c|c|} \hline
$E_{1}$ & $E_{2}$ & $E_{1}$ {\idt ||} $E_{2}$\\ \hline \hline
{\idt true} & {\idt true} & {\idt true}\\ \hline
{\idt true} & {\idt false} & {\idt true}\\ \hline
{\idt false} & {\idt true} & {\idt true}\\ \hline
{\idt false} & {\idt false} & {\idt false}\\ \hline
\end{tabular}
\end{center}

La conjonction et la disjonction peuvent se réaliser simplement avec un circuit électrique comme on le voit à la figure
\ref{fig:andor}
ci-dessous
%(page \pageref{fig:andor})
avec la con\-ven\-tion que {\idt true} représente l'interrupteur fermé
et l'ampoule allumée (le courant passe) et {\idt false} représente l'interrupteur ouvert
et l'ampoule éteinte (le courant ne passe pas).

\begin{figure}[htbp]
 \includevspacePostScript[4.3cm]{CircuitsAndOr}{0.95}{0.95}
 \caption{Circuits électriques réalisant la conjonction et\label{fig:andor}}
\centering{la disjonction logiques}
\end{figure}

Une disjonction est utilisée dans le programme suivant pour tester si le \robot\ est orienté au nord ou à
l'est~:

\vspace{0.3cm}
 {\raggedright
 \includevspacePostScript[0.5cm]{FriseEnOr}{1.0}{1.0}
{
\parindent \parindentprog
\medskip

\carrenoir
\vspace{\neglskip}\mbox{\blc {\mcl open}{\idt \ LibRobot}}\par
\vspace{0.3\baselineskip}
\vspace{\neglskip}\mbox{\blc {\cmt (*\ Dessiner\ une\ frise\ en\ or\ *)}}\par
\vspace{0.3\baselineskip}
\vspace{\neglskip}\mbox{\blc {\mcl let}{\idt \ longueur\ =\ 31}}\par
\vspace{0.3\baselineskip}
\vspace{\neglskip}\mbox{\blc {\mcl let}{\idt \ ()\ =}}\par
\vspace{\neglskip}\mbox{\blc \blk{\idt ag\ ();}}\par
\vspace{\neglskip}\mbox{\blc \blk{\mcl for}{\idt \ i\ =\ 1\ }{\mcl to}{\idt \ longueur\ }{\mcl do}}\par
\vspace{\neglskip}\mbox{\blc \blk\blk{\mcl if}{\idt \ tn\ ()\ \char124{}\char124{}\ te\ ()\ }{\mcl then}}\par
\vspace{\neglskip}\mbox{\blc \blk\blk\blk{\idt vd\ ()}}\par
\vspace{\neglskip}\mbox{\blc \blk\blk{\mcl else}{\idt \ }{\mcl begin}}\par
\vspace{\neglskip}\mbox{\blc \blk\blk\blk{\idt avf\ ({-}1);\ vg\ ();\ pqtg\ ()}}\par
\vspace{\neglskip}\mbox{\blc \blk\blk{\mcl end}}\par
\vspace{\neglskip}\mbox{\blc \blk{\mcl done}{\idt ;}}\par
\vspace{\neglskip}\mbox{\blc \blk{\idt st\ ()}}\par

\medskip }

  }

\section{Expressions booléennes}

La table ci-dessous résume les règles d'écriture et les règles de calcul de la valeur
des expressions booléennes en \ocaml~:

\index{Expression booléenne} \index{Test}

\begin{center}
\protect\frenchspacing
\begin{tabular}{@{}|l|p{6.8cm}|@{}} \hline
\multicolumn{2}{|l|}{\em Expressions booléennes~:} \\ \hline \hline
{\idt true} & Valeur `vrai'.\\
{\idt false} & Valeur `faux'.\\ \hline
{\small\rm Expr. entière$_{1}$} {\idt <=} {\small\rm Expr. entière$_{2}$}& {\idt true} si et seulement si la valeur de l'{\small\rm Expr. entière$_{1}$} est
                                 inférieure ou égale à celle de l'{\small\rm Expr. entière$_{2}$}.\\
{\small\rm Expr. entière$_{1}$} {\idt <} {\small\rm Expr. entière$_{2}$}& {\idt true} si et seu\-le\-ment si la va\-leur de l'{\small\rm Expr. entière$_{1}$} est
                                 strictement inférieure à celle de l'{\small\rm Expr. entière$_{2}$}.\\
{\small\rm Expr. entière$_{1}$} {\idt =} {\small\rm Expr. entière$_{2}$}& {\idt true} si et seu\-le\-ment si la valeur de l'{\small\rm Expr. entière$_{1}$} est
                                 égale à celle de l'{\small\rm Expr. entière$_{2}$}.\\
{\small\rm Expr. entière$_{1}$} {\idt <>} {\small\rm Expr. entière$_{2}$}& {\idt true} si et seu\-le\-ment si la valeur de l'{\small\rm Expr. entière$_{1}$} est
                                 différente de celle de l'{\small\rm Expr. entière$_{2}$}.\\ \hline
\end{tabular}
\end{center}
\begin{center}
\protect\frenchspacing
\begin{tabular}{@{}|l|p{6.8cm}|@{}} \hline
{\small\rm Expr. entière$_{1}$} {\idt >} {\small\rm Expr. entière$_{2}$}& {\idt true} si et seu\-le\-ment si la va\-leur de l'{\small\rm Expr. entière$_{1}$} est
                                 strictement supérieure à celle de l'{\small\rm Expr. entière$_{2}$}.\\
{\small\rm Expr. entière$_{1}$} {\idt >=} {\small\rm Expr. entière$_{2}$}& {\idt true} si et seulement si la valeur de l'{\small\rm Expr. entière$_{1}$} est
                                supérieure ou égale à celle de l'{\small\rm Expr. entière$_{2}$}.\\ \hline
                                \end{tabular}
\end{center}
\begin{center}
\protect\frenchspacing
\begin{tabular}{@{}|l|p{6.8cm}|@{}} \hline
{\mcl not} {\idt (}{\small\rm Expr. bool.}{\idt )} & {\idt true} si et seulement si la valeur de l'{\small\rm Expr. bool.} est
                                égale à {\idt false}.\\ \hline
{\small\rm Expr. bool.$_{1}$} {\idt \&\&} {\small\rm Expr. bool.$_{2}$}& {\idt true} si et seulement si les valeurs de l'{\small\rm Expr. bool.$_{1}$} et
                                 de l'{\small\rm Expr. bool.$_{2}$} sont égales à {\idt true}.\\
{\small\rm Expr. bool.$_{1}$} {\idt ||} {\small\rm Expr. bool.$_{2}$}& {\idt true} si et seulement si la valeur de l'{\small\rm Expr. bool.$_{1}$}
                                 ou celle de l'{\small\rm Expr. bool.$_{2}$} est égale à {\idt true}.\\ \hline
\end{tabular}
\end{center}

\index{Fonction booléenne}On peut donner un nom à une expression booléenne pour
rendre les programmes plus lisibles. Par exemple~:

{
\parindent \parindentprog
\medskip

\vspace{\neglskip}\mbox{\blc {\cmt (*\ n\ est-il\ impair\ ?\ *)}}\par
\vspace{\neglskip}\mbox{\blc {\mcl let}{\idt \ estImpair\ n\ =\ n\ }{\mcl mod}{\idt \ 2\ =\ 1}}\par
\vspace{0.3\baselineskip}
\vspace{\neglskip}\mbox{\blc {\cmt (*\ n\ est-il\ pair\ ?\ *)}}\par
\vspace{\neglskip}\mbox{\blc {\mcl let}{\idt \ estPair\ n\ =\ n\ }{\mcl mod}{\idt \ 2\ =\ 0}}\par

\medskip }

\noindent Ces déclarations ressemblent à des déclarations de procédures avec paramètres
mais ce qui suit le signe {\idt =} n'est pas une suite d'instructions pour le
\robot~: c'est une expression booléenne dont la valeur est calculée à chaque appel.
On dit que {\idt estImpair} et {\idt estPair} sont des {\em fonctions booléennes} car leur évaluation retourne toujours un booléen. Comme leur paramètre est entier, on dit qu'elles ont pout type \texttt{int $\rightarrow$ bool}.
Par exemple {\idt estImpair 5} vaut {\idt true} et {\idt estPair 5} vaut
{\idt false}. On peut les utiliser dans un test comme {\mcl if}\ {\idt estImpair i}\ {\mcl then}
\ldots\ ou dans une expression booléenne comme {\mcl not}\ {\idt (estImpair i)}. Les tests
d'orientation {\idt tn ()}, {\idt te ()}, \ldots\ du \robot\ sont également des fonctions
booléennes, déclarées dans la bibliothèque {\idt LibRobot}.

\noindent \rule{137mm}{0.1mm} %---------------------------------------------
\begin{exercices}{Écrire des programmes \ocaml\  utilisant des tests de parité pour réaliser
les dessins suivants~:

\tiret  la serpentine pyramidale et la grille de la figure \ref{fig:serpyrgricr} ci-dessous$\:$;
%(page \pageref{fig:serpyrgricr})$\:$;

\tiret   le pavage octogonal et le pavage d'hexagones entrelacés de la figure \ref{fig:pavocthex} (page \pageref{fig:pavocthex}).

\noindent $\Box$}\end{exercices}
\begin{figure}[htbp]
  \includevspacePostScript[6.0cm]{SerpPyr+GrilleDeCr}{0.95}{0.95}
 \caption{Serpentine pyramidale et grille de croix}\label{fig:serpyrgricr}
\end{figure}
\begin{figure}[htbp]
\includevspacePostScript[5.4cm]{PavageOctoHexa}{0.95}{0.95}
 \caption{Pavages d'octogones juxtaposés et d'hexagones entrelacés}\label{fig:pavocthex}
\end{figure}
\begin{exercices}{Écrire des programmes \ocaml\  utilisant des expressions boolé\-en\-nes de comparaison de
valeurs entières pour réaliser les dessins suivants (le deuxième est difficile)~:

\tiret  les chiffres digitaux de la figure \ref{fig:alphabet} (page \pageref{fig:alphabet}) où le chiffre à dessiner
est passé en paramètre de la procédure de dessin$\:$;

\tiret  le pavage de roues dentées de la figure \ref{fig:pavrd} ci-dessous$\:$;

\tiret  le tapis égyptien de la figure \ref{fig:tapegyp} (page \pageref{fig:tapegyp}).

\noindent $\Box$}\end{exercices}
\begin{figure}[htbp]
 \includevspacePostScript[5.7cm]{PavageRoueDentee}{0.95}{0.95}
 \caption{Pavage de roues dentées}\label{fig:pavrd}
\end{figure}
\begin{figure}[htbp]
 \includevspacePostScript[7.2cm]{TapisEgyptien}{0.95}{0.95}
 \caption{Tapis égyptien\label{fig:tapegyp}}
\end{figure}
\begin{exercices}{Écrire des programmes \ocaml\  utilisant des négations, conjonctions ou disjonctions d'expressions
booléen\-nes pour réaliser les dessins suivants (le deuxième est difficile)~:

\tiret  le carré zébré de la figure \ref{fig:carzeb} (page \pageref{fig:carzeb})$\:$;

\tiret  la mosaïque de l'Alhambra représentée à la figure \ref{fig:mosaic6} (page \pageref{fig:mosaic6}).

\noindent $\Box$}\end{exercices}
\begin{figure}[htbp]
 \includevspacePostScript[5.8cm]{CarreZebree}{0.95}{0.95}
 \caption{Carré zébré}\label{fig:carzeb}
\end{figure}
\begin{figure}[htbp]
 \includevspacePostScript[5.7cm]{MosaiqueAlhambra6}{0.95}{0.95}
 \caption{Mosaïque de l'Alhambra}\label{fig:mosaic6}
\end{figure}
%%%%%%%%%%%%%%%%%%%%%%%%%%%%%%%%%%%%%%%%%%%%%%%%%%%%%%%%%%
\begin{corrige}{
Nous donnons un programme, les autres se trouvant dans le dossier
{\idt livre/\allowbreak 16\_expressions\_booleennes}~:
\input ocaml/SerpentinePyramidale.tex
\noindent $\Box$}\end{corrige}
\begin{corrige}{
Nous donnons un programme, le troisième se trouvant dans le dossier
{\idt livre/\allowbreak 16\_expressions\_booleennes}~:
\input ocaml/PavageRouesDentees.tex
\noindent $\Box$}\end{corrige}
\begin{corrige}{
Nous donnons un programme, le premier se trouvant dans le dossier
{\idt livre/\allowbreak 16\_expressions\_booleennes}~:
\input ocaml/MosaiqueAlhambra6.tex
\noindent $\Box$}\end{corrige}

%%%%%%%%%%%%%%%%%%%%%%%%%%%%%%%%%%%%%%%%%%%%%%%%%%%%%%%%%%%%%%%%%%%%%%%%%%%
\noindent \framebox[\textwidth]{
\begin{minipage}{\myparboxwidth}
\setlength{\parindent}{2em}
\small \sf
\vspace{\abovedisplayskip}
\section*{Addenda}

 En \ocaml\ les priorités des opérateurs entiers et booléens sont définies comme suit~:

 \vspace{\baselineskip}

 \index{Priorité des opérateurs}
\begin{center}
\begin{tabular}{|l|l|l|} \hline
\multicolumn{3}{|l|}{\em Priorité des opérateurs entiers et booléens~:} \\ \hline \hline
{\idt f x} (appel), $-$ (unaire)  & appels de fonctions & priorité la plus forte \\ \hline
$\ast$, {\idt /}, {\mcl mod} & opérateurs multiplicatifs & \\ \hline
$+$, $-$ (binaires) & opérateurs additifs & \\ \hline
{\idt <}, {\idt <=}, {\idt =}, {\idt <>}, {\idt >=}, {\idt >}  & opérateurs de relation & \\ \hline
{\idt \&\&} & conjonction & \\ \hline
{\idt ||} & disjonction & priorité la plus faible \\ \hline
\end{tabular}
\end{center}

\vspace{\baselineskip}

 \noindent Par exemple l'expression booléenne `{\idt  true\ \&\&\ false\ ||\ false}' s'é\-va\-lue com\-me
 `{\idt  (true\ \&\&\ false)\ ||\ false}' et est donc égale à {\idt false}.

\vspace{\abovedisplayskip}
\end{minipage}
}
%%%%%%%%%%%%%%%%%%%%%%%%%%%%%%%%%%%%%%%%%%%%%%%%%%%%%%%%%%%%%%%%%%%%%%%%%%%%

%%%%%%%%%%%%%%%%%%%%%%%%%%%%%%%%%%%%%%%%%%%%%%%%%%%%%%%%%%%%%%%%%%%%%%%%%%%
\noindent \framebox[\textwidth]{
\begin{minipage}{\myparboxwidth}
\setlength{\parindent}{2em}
\small \sf
\vspace{\abovedisplayskip}
\indent Ces priorités sont proches des conventions mathématiques usuelles. Par
exemple `{\idt  0\ <=\ i\ \&\&\ i\ <\ 100}' s'évalue comme
`{\idt  (0\ <=\ i)\ \&\&\ (i\ <\ 100)}', ce qui est bien ce que l'on veut dire. Il n'en était
pas de même dans des langages plus anciens comme \pascal, où cette expression était
incorrecte et où il fallait écrire toutes les parenthèses.

Comme {\mcl not} s'applique à ce qui le suit immédiatement, à la manière d'un appel de
fonction, `{\mcl not}\ {\idt i\ mod\ 2\ =\ 1}' s'évalue comme `{\idt (}{\mcl not}\ {\idt i)\ mod\ 2\ =\ 1}'
qui est refusé par le compilateur puisque `{\mcl not}\ {\idt i}' est un booléen, que l'on ne peut pas diviser par 2. Il faut
donc écrire `{\mcl not}\ {\idt (i\ mod\ 2\ =\ 1)}'. Dans le doute, on peut toujours
ajouter des parenthèses.

Enfin, \ocaml\ évalue les conjonctions et les disjonctions de gauche à droite et s'arrête dès
que le résultat est connu~: dans $E_{1}$ {\idt \&\&} $E_{2}$, si $E_{1}$ vaut {\idt false}
alors $E_{2}$ n'est pas calculée puisque la conjonction vaut forcément {\idt false}$\:$; de
même dans $E_{1}$ {\idt ||} $E_{2}$, si $E_{1}$ vaut {\idt true} alors $E_{2}$ n'est pas
calculée.

Le terme ``booléen'' ou ``booléien'' ({\idt bool} en \ocaml) vient de Boole\footnote{\ \href{https://fr.wikipedia.org/wiki/George_Boole}{George Boole}, mathématicien britannique, 1815--1864.},
 créateur de la logique mathématique moderne.

\vspace{\abovedisplayskip}
\end{minipage}
}
%%%%%%%%%%%%%%%%%%%%%%%%%%%%%%%%%%%%%%%%%%%%%%%%%%%%%%%%%%%%%%%%%%%%%%%%%%%%
\newpage                 % Si on termine le chapitre par une page blanche
\thispagestyle{empty}
